% ============================================================
% Introduction — README §1
%   §1.1 Provenance-aware memory is a hot topic, bound to honest writers
%   §1.2 But the attacker is the writer
%   §1.3 Attribution as behavioral trace, not field
%   §1.4 Headline: In-Record Attribution Verification (R3)
%   §1.5 Contributions
% ============================================================
\section{Introduction}
\label{sec:intro}

% ---------- §1.1 honest-writer premise ----------
\paragraph{Provenance-aware memory is bound to honest writers.}
\todo{Open with the active line of work (\tiermem, MemOS MemCube,
enterprise stamping) and the structural assumption they share:
metadata fields are facts because the writer is trusted. Cite
\tiermem{} and MemOS preprints. End with the failure mode: once
memory leaves the original system, that assumption collapses.}

% ---------- §1.2 attacker is the writer ----------
\paragraph{The attacker is the writer.}
\todo{Enumerate what the snapshot-holder can do: stamp ownership,
erase \texttt{source\_agent} / \texttt{tenant\_id}, fabricate
provenance chains, rewrite the memory system's own changelog
(\amem{} evolve history, \graphiti{} fact-invalidation chain) since
they sit in the same store. Conclude: any field-based provenance is
ill-defined in adversarial storage. Token-level LLM watermarks and
agent action-layer watermarks don't fill the gap — they assume the
trajectory survives, which it doesn't in the most forensically
valuable case.}

% ---------- §1.3 attribution = behavioral trace ----------
\paragraph{Attribution as behavioral trace, not field.}
\todo{Frame the three carriers (update / link / semantic) as latent
decisions every memory write makes. Properties: (i) non-trivial
candidate space, (ii) equivalent choices don't break utility, (iii)
the choice itself is readable from the snapshot. Then embed a PRF
key into the sequence of keyed picks: $\nonceT = \text{PRF}(K, \ctx)$
selects which equivalent candidate is realized. Reads as ``\emph{the
writer's key can be replayed against the snapshot}'' rather than
``\emph{we say so}''.}

% ---------- §1.4 R3 headline ----------
\paragraph{In-Record Attribution Verification (R3).}
\todo{Define the three regimes R1 (full external log) / R2 (partial
log) / R3 (snapshot only). R3 is the headline: bind each keyed
selection into a per-decision commitment, anchor the per-session
Merkle root + signature directly inside the snapshot, and store the
reveal record co-located with the memory record itself. Forward to
\S\ref{sec:rq3} for the numbers.}

% ---------- §1.5 contributions (3 pillars) ----------
% Three threads of the paper:
%   (i)   The watermark mechanism (carrier taxonomy + interception
%         + keyed sampling + cryptographic audit trace, packaged as
%         one engineering contribution).
%   (ii)  In-Record Attribution Verification (R3), the headline
%         result that no prior watermark or provenance scheme
%         provides under adversarial storage.
%   (iii) Memory-specific evaluation showing utility preservation
%         within statistical noise + robustness under 9 attacks.
\paragraph{Contributions.}
\begin{enumerate}[topsep=2pt,itemsep=2pt,leftmargin=*]
\item \textbf{A watermark for long-term agent memory.}
  \todo{One sentence on the end-to-end mechanism: backend-invariant
  $\langle \carrier, \candset, \probdist, \ctx \rangle$ taxonomy
  (\S\ref{sec:method-formal}), \agentmark-style 1-call enumeration
  that produces $(\candset, \probdist)$ from open-vocabulary
  structured-JSON LLM output (\S\ref{sec:method-interception}),
  distribution-preserving keyed sampling on the SDK's internal LLM
  calls (\S\ref{sec:method-sampler}), and a cryptographic audit
  trace (commitment + per-session Merkle log + signed in-snapshot
  anchor + per-leaf inclusion proof; \S\ref{sec:method-audit})
  that supports snapshot-only verification.}
\item \textbf{In-Record Attribution Verification (R3).}
  \todo{One sentence on the three-regime hierarchy and the
  empirical result: \memmark{} recovers the full payload from
  snapshot alone with bit\_recovery $= 1.00$ on \locomo{}, while
  a \texttt{signed-metadata-only} baseline drops to $0$ and the
  wrong-key control drops to $0.15$ (\S\ref{sec:rq3}). No prior
  watermark or provenance scheme provides snapshot-only attribution
  under the adversarial-storage threat model.}
\item \textbf{Memory-specific evaluation across two backends and
  two LLMs.} \todo{One sentence: utility preserved within
  statistical noise on \locomo{} (\S\ref{sec:rq1}); non-trivial
  capacity (\S\ref{sec:rq2}); robust under nine memory-lifecycle
  attacks including KGMark-style edge relabel and subgraph
  reanchor (\S\ref{sec:rq4}); ingestion integrity unaffected
  (\S\ref{sec:rq5}).}
\end{enumerate}

% ---------- pointer figure ----------
% ============================================================
% Figure: Overview of MemMark's storyline (3-panel)
%   panel A: provenance-as-field (attacker forges)
%   panel B: keyed-decision behavioral trace (attacker can't forge)
%   panel C: R3 in-record verification (snapshot-only)
% ============================================================
\begin{figure*}[t]
\centering
\begin{tikzpicture}[
  node distance=8mm and 14mm,
  panel/.style={draw, thick, rounded corners=2pt, minimum width=0.30\textwidth, minimum height=4.2cm, align=center, inner sep=4pt},
  attacker/.style={fill=red!8},
  honest/.style={fill=green!5},
  reader/.style={fill=blue!5},
  title/.style={font=\bfseries\small},
]

% Panel A — honest writer assumption
\node[panel, attacker] (A) {
  \begin{minipage}{0.28\textwidth}
  \centering
  \todo{\textbf{Panel A — Honest-Writer Provenance.}\\
  Metadata fields claim authorship. Attacker who holds the
  snapshot rewrites any field. Token-level / action-layer
  watermarks need trajectory, which is gone.}
  \end{minipage}
};

% Panel B — keyed behavioral trace
\node[panel, honest, right=of A] (B) {
  \begin{minipage}{0.28\textwidth}
  \centering
  \todo{\textbf{Panel B — Keyed Behavioral Trace.}\\
  Each evolve decision $\decision$ is a non-trivial choice over
  $\candset$. PRF key $K$ selects which one is realized via
  binning sampler. Choice itself is what gets stored.}
  \end{minipage}
};

% Panel C — R3 in-record verification
\node[panel, reader, right=of B] (C) {
  \begin{minipage}{0.28\textwidth}
  \centering
  \todo{\textbf{Panel C — In-Record Verification (R3).}\\
  Per-decision commitment + per-session Merkle anchor + per-leaf
  inclusion proof live inside the snapshot. Verifier with key
  $K$ alone recovers payload bits.}
  \end{minipage}
};

% Title labels above
\node[title, above=1mm of A] {\textsc{Provenance as Field}};
\node[title, above=1mm of B] {\textsc{Provenance as Behavior}};
\node[title, above=1mm of C] {\textsc{Snapshot-Only Attribution}};

\end{tikzpicture}
\caption{\textbf{\memmark{} in three panels.} (A) Field-based
provenance assumes the writer is honest; once the attacker holds
the snapshot, all fields are writable. (B) \memmark{} embeds
attribution into the memory system's own latent evolve decisions
via a PRF-keyed binning sampler that preserves the marginal
$\probdist$ distribution per Lemma~\ref{lem:dist-preserve}. (C) The
in-record audit trace (commitment + Merkle anchor + per-leaf
inclusion proof) lets a verifier with only the snapshot and key
$K$ decode the writer's bits — RQ3 headline.}
\label{fig:overview}
\end{figure*}

